\documentclass[10pt,a4paper]{article}%
\usepackage[T1]{fontenc}%
\usepackage[utf8]{inputenc}%
\usepackage{lmodern}%
\usepackage{textcomp}%
\usepackage{lastpage}%
\usepackage{geometry}%
\geometry{margin=0.65in}%
%
\usepackage{enumitem}%
\usepackage{hyperref}%
\usepackage{titlesec}%
\usepackage{xcolor}%
\titleformat{\section}{\large\bfseries\uppercase}{}{0em}{}[\titlerule]%
\titlespacing{\section}{0pt}{10pt}{6pt}%
\pagestyle{empty}%
%
\begin{document}%
\normalsize%
\noindent {\fontsize{26}{32}\selectfont \textbf{Sanshray Choudhary}} \vspace{3pt}\\%
\noindent \textbf{\scriptsize sanshray.choudhary@gmail.com $\bullet$ +91 9810894391 $\bullet$ \href{https://www.linkedin.com/in/sanshray-choudhary-3591b9282}{https://www.linkedin.com/in/sanshray-choudhary-3591b9282} \\ \href{https://github.com/sanshraychoudhary-hue}{https://github.com/sanshraychoudhary-hue}} \vspace{6pt}%

\section*{Target Position: Robotics Software Intern -- Antropi Robotics (YC F26)}%
\noindent \textit{Electronics \& Communication Engineer (USICT, GPA 8.42) specializing in ARM Cortex/STM32, ESP32-S3 FreeRTOS firmware, custom KiCad PCB hardware design, and protocol integration (CAN, Modbus, I2S, SPI, I2C). e-Yantra IIT Bombay Certified in Embedded Systems \& ROS 2.} \vspace{6pt}%

\section*{Education}%
\noindent \textbf{University School of Information, Communication and Technology, Dwarka - 110078} \hfill 2023--2027 \\%
\indent \textit{Bachelor of Technology, Electronics and Communication Engineering (8.42 GPA)} \vspace{6pt}\\%
\noindent \textbf{The Air Force School, Subroto Park, Delhi Cantt - 110010} \hfill 2009--2022 \\%
\indent \textit{Senior Secondary (XII) (PCM), CBSE (93.8\%) $\bullet$ Secondary (X), CBSE (94.6\%)} \vspace{6pt}%

\section*{Licenses \& Certifications}%
\noindent \textbf{e-Yantra Robotics Competition (Level 3)} \hfill Mar 2025 \\%
\indent \textit{Issued by e-Yantra, IIT Bombay} \vspace{4pt}\\%
\noindent \textbf{ROS 2 Developer} \hfill Feb 2025 \\%
\indent \textit{Issued by e-Yantra, IIT Bombay} \vspace{4pt}\\%
\noindent \textbf{Embedded Systems Developer} \hfill Dec 2024 \\%
\indent \textit{Issued by e-Yantra, IIT Bombay} \vspace{6pt}%

\section*{Professional Experience}%
\noindent \textbf{Winter Intern, India Space Academy} \hfill Feb 2026 -- Mar 2026%
\begin{itemize}[leftmargin=15pt, nosep, topsep=0pt]%
\item Developed embedded firmware for advanced drone technology utilizing microcontrollers, I2C/SPI interfaces, and optimized power management.%
\end{itemize} \vspace{6pt}%

\noindent \textbf{Project Intern, AICTE (OmniServe Robot)} \hfill Jun 2025 -- Aug 2025%
\begin{itemize}[leftmargin=15pt, nosep, topsep=0pt]%
\item Engineered an autonomous Wi-Fi mobile robot using Raspberry Pi Pico W, mecanum wheels, and closed-loop encoder PID control.%
\end{itemize} \vspace{6pt}%

\noindent \textbf{TechSpace, Joint Secretary} \hfill Nov 2024 -- Present%
\begin{itemize}[leftmargin=15pt, nosep, topsep=0pt]%
\item Facilitated technical events, workshops, and embedded systems seminars to foster peer-to-peer technical growth.%
\end{itemize} \vspace{6pt}%

\noindent \textbf{Internship Trainee, BSNL} \hfill Jun 2026 -- Jul 2026 \vspace{6pt}%

\section*{Highlighted Relevant Projects}%
\noindent \textbf{OmniServe (Internship Project)} \hfill Jun 2025%
\begin{itemize}[leftmargin=15pt, nosep, topsep=0pt]%
\item Developed an autonomous Wi-Fi controlled omni-directional mobile robot using Raspberry Pi Pico W with mecanum wheels, encoder-based PID motion control, and centimeter-level navigation accuracy.%
\item Implemented a web-based control dashboard with real-time telemetry and integrated a robotic arm for basic pick-and-place operations, validated under real-world load conditions.%
\end{itemize} \vspace{8pt}

\noindent \textbf{Self-Balancing Bot - eYantra 2024 (Stage 1 – Qualified)} \hfill Sep 2024%
\begin{itemize}[leftmargin=15pt, nosep, topsep=0pt]%
\item Built a self-balancing robot using an Arduino Nano, MPU6050, and encoders. Tuned the PID controller and simulated in Python/CoppeliaSim.%
\end{itemize} \vspace{8pt}

\noindent \textbf{AtiDhwani -- ESP32-S3 Ultrasonic Bat PAM Node \& Bioacoustics Detector} \hfill Nov 2026 -- Present%
\begin{itemize}[leftmargin=15pt, nosep, topsep=0pt]%
\item Engineered an open-source modular Passive Acoustic Monitoring (PAM) node and handheld bat detector capturing 10--96\,kHz wildlife bioacoustics with 192\,kHz / 24-bit continuous sampling at $<$ \$40 BOM cost (vs.\ \$1,000--\$3,500 commercial systems).%
\item Developed dual-core FreeRTOS firmware on ESP32-S3 (240\,MHz, 8MB Octal PSRAM): Core 0 manages high-speed DMA I2S streaming and drop-free FAT32 WAV logging to MicroSD via a 3.6MB circular PSRAM buffer.%
\item Implemented real-time DSP on Core 1 executing quadrature heterodyne down-conversion (tunable 15--90\,kHz LO), frequency division, and 128-point Radix-4 FFT spectral analysis displayed on a 0.91" SSD1306 OLED at 30\,FPS.%
\item Designed custom 2-layer carrier PCB in KiCad (0 DRC violations) integrating dual MEMS microphones, ultra-low-noise LDO, 24-bit Wolfson WM8782 stereo ADC (Master Mode, 24.576\,MHz crystal), and 32-bit TI PCM5102A DAC driving field headphones with zero perceived latency.%
\end{itemize} \vspace{8pt}

\noindent \textbf{STM32 Industrial IoT Sensor Hub} \hfill Jan 2025%
\begin{itemize}[leftmargin=15pt, nosep, topsep=0pt]%
\item Developed a high-performance industrial sensor node using STM32F407, interfacing with SPI-based ADCs, I2C temperature sensors, and UART gas sensors.%
\item Implemented a lightweight FreeRTOS firmware to manage concurrent sensor sampling, data filtering, and reliable MQTT transmission over an ESP8266 Wi-Fi module.%
\item Designed an intelligent power management algorithm using deep sleep modes to maximize battery life in remote deployments.%
\end{itemize} \vspace{8pt}

\noindent \textbf{Aegis Guardian - AI Fall Detection \& Health Monitoring System} \hfill Oct 2025%
\begin{itemize}[leftmargin=15pt, nosep, topsep=0pt]%
\item Engineered an ESP32-based embedded wearable system for real-time fall detection and vital monitoring.%
\item Developed firmware for 50Hz sensor acquisition (MPU6050, AD8232, BMP180) using I2C/ADC with optimized Wi-Fi data streaming to a Python ML inference server.%
\item Implemented a multi-model pipeline (Random Forest + TFLite CNN) achieving 99.97\% accuracy with sub-50ms latency, and integrated cloud sync for remote monitoring.%
\end{itemize} \vspace{8pt}

\section*{Core Technical Skills}%
\noindent \textbf{Core Electronics \& Embedded Systems:} Microcontrollers (ARM Cortex M3/M4, STM32, ESP32, Arduino), IoT, Robotics, Signal Processing (FFT, Filtering), Digital Electronics, Circuit Design, PCB Design (KiCad), FreeRTOS, Sensors \& Actuators, Control Systems, PID Control, Power Electronics, CNC Machining, STM32Cube, Protocols \& Interfaces (SPI, I2C, I2S, UART, CAN, RS-485, Modbus, MQTT), Hardware Debugging (Logic Analyzers, Oscilloscopes) \vspace{6pt}\\%
\noindent \textbf{Software \& Tools:} C, C++, Python, MATLAB, TinyML, TensorFlow Lite, OpenCV, ROS/ROS 2, Git, Linux, Simulation \& Modelling (CoppeliaSim)%
\end{document}